% Garder le prélèvement forfaitaire unique (rapport Sorbon, partie III)
% Compilation depuis fiscalite/pfu/ : pdflatex note-pfu.tex (deux passes)
% Matière et statut des sources : notes/pfu_recherche.md ; chiffres : donnees/pfu_chiffres.csv
% Figure : python3 pfu/figures_pfu.py puis FIG_NU=1 python3 pfu/figures_pfu.py (depuis fiscalite/)
\documentclass[11pt,a4paper]{article}
% Préambule commun au brief et à la revue de littérature sur la taxe sur la valeur foncière.
% Même composition que le cours (cours/refaire-l-impot.tex) ; figures « nues » produites par
%   FIG_NU=1 python lvt/graphiques_lvt.py
\usepackage[T1]{fontenc}
\usepackage[utf8]{inputenc}
\usepackage{lmodern}
\usepackage[french]{babel}
\usepackage[a4paper,left=2.8cm,right=2.8cm,top=2.6cm,bottom=2.8cm]{geometry}
\usepackage{graphicx}
\usepackage[export]{adjustbox}
\usepackage{booktabs}
\usepackage{tabularx}
\usepackage{array}
\usepackage{xcolor}
\usepackage{titlesec}
\usepackage[labelfont=bf,labelsep=period,font=small,singlelinecheck=false,
            justification=raggedright,skip=6pt]{caption}
\usepackage{fancyhdr}
\usepackage{enumitem}
\usepackage{microtype}
\usepackage{textcomp}
\usepackage{amsmath}
\providecommand{\euro}{\texteuro}
\usepackage[hidelinks]{hyperref}

\definecolor{encre}{HTML}{1A1A19}
\definecolor{bleu}{HTML}{003399}
\definecolor{vert}{HTML}{1F7A4D}
\definecolor{ocre}{HTML}{C79100}
\definecolor{brique}{HTML}{993333}
\definecolor{filet}{HTML}{C8C6BE}

\titleformat{\section}{\normalfont\large\bfseries\color{bleu}}{\thesection.}{0.6em}{}
\titleformat{\subsection}{\normalfont\normalsize\bfseries\color{bleu}}{\thesubsection}{0.6em}{}
\titlespacing*{\section}{0pt}{2.4ex plus 1ex minus .2ex}{1.2ex plus .2ex}
\titlespacing*{\subsection}{0pt}{2ex plus 1ex minus .2ex}{0.8ex plus .2ex}

\pagestyle{fancy}
\fancyhf{}
\renewcommand{\headrulewidth}{0.4pt}
\fancyhead[R]{\small\thepage}

\setlength{\parindent}{0pt}
\setlength{\parskip}{0.62em}
\linespread{1.05}
\renewcommand{\arraystretch}{1.2}
\setlist{itemsep=2pt,topsep=2pt,parsep=0pt,leftmargin=1.4em}

\newenvironment{encadre}[1]
  {\par\medskip\noindent\begin{minipage}{\linewidth}%
   \color{encre}\setlength{\parskip}{0.5em}%
   \hrule height 0.8pt width \linewidth \vspace{6pt}%
   {\footnotesize\bfseries\color{vert}\MakeUppercase{#1}}\par\vspace{-2pt}%
   \small}
  {\par\vspace{4pt}\hrule height 0.4pt width \linewidth\end{minipage}\par\medskip}

\newcommand{\figbloc}[3]{%
  \begin{figure}[htbp]\centering
  \caption{#2}\label{#3}
  \includegraphics[max width=\linewidth,max height=0.40\textheight]{../figures/nu/#1.pdf}
  \IfFileExists{../figures/nu/#1_note.tex}{\par\vspace{4pt}%
    \parbox{\linewidth}{\raggedright\scriptsize\color{encre}\setlength{\parskip}{2pt}%
    \input{../figures/nu/#1_note.tex}}}{}%
  \end{figure}}

\newcommand{\pc}{\,\%}

\renewcommand{\figbloc}[3]{%
  \begin{figure}[!htbp]\centering
  \caption{#2}\label{#3}
  \includegraphics[max width=\linewidth,max height=0.27\textheight]{../figures/nu/#1.pdf}
  \par\vspace{3pt}\parbox{\linewidth}{\raggedright\scriptsize\color{encre}\setlength{\parskip}{2pt}%
    \input{../figures/nu/#1_note.tex}}%
  \end{figure}}
\usepackage{needspace}
\usepackage{etoolbox}
\preto\section{\needspace{6\baselineskip}}
\renewcommand{\topfraction}{0.9}\renewcommand{\bottomfraction}{0.8}
\renewcommand{\textfraction}{0.08}\renewcommand{\floatpagefraction}{0.85}
\fancyhead[L]{\small\itshape Garder le prélèvement forfaitaire unique}

\begin{document}
\renewcommand{\figurename}{Figure}\renewcommand{\tablename}{Tableau}

\thispagestyle{empty}
\vspace*{0.6cm}
{\footnotesize\color{ocre}\bfseries NOTE DE TRAVAIL\quad·\quad NOVEMBRE 2026\par}
\vspace{8pt}
{\fontsize{22}{26}\selectfont\bfseries\color{bleu}Garder le prélèvement forfaitaire unique\par}
\vspace{6pt}
{\large Un taux qui a fait ses preuves ; la faille est dans le circuit des holdings\par}
\vspace{8pt}
{\small Colin Baget. Droit en vigueur au 15 novembre 2026 (LFSS et LF 2026) ; évaluations de France
Stratégie, de l'IPP et de l'Insee ; rapport de la commission d'enquête de l'Assemblée nationale n°~3056.
Sources et statut de lecture : \texttt{fiscalite/pfu/notes/}.\par}
\vspace{6pt}
{\color{bleu}\hrule height 1.2pt}
\vspace{4pt}

\section*{Introduction}

Le prélèvement forfaitaire unique (PFU) impose depuis 2018 les dividendes, les intérêts et les plus-values
de valeurs mobilières à un taux unique, de 31,4~\% en 2026, au lieu du barème progressif. Cette note
recommande de le garder tel qu'il est : même taux, même structure. Elle écarte les deux idées qui
circulent pour le modifier, une baisse à 25~\% et un impôt de 25~\% sur les revenus conservés dans les
holdings. Elle retient une seule correction, déjà inscrite au projet de loi de finances pour 2027, et
précise comment le PFU s'articule avec la fusion de l'impôt sur le revenu et de la CSG.

La raison principale est que la France a fait l'expérience dans les deux sens. En 2013, l'imposition des
dividendes au barème a fait chuter de 40~\% les dividendes déclarés par les foyers concernés, et la réforme
a coûté de l'argent au lieu d'en rapporter. En 2018, le PFU les a fait remonter et a coûté beaucoup moins que
prévu. Le PFU ne doit pas pour autant être présenté comme une réforme favorable à l'investissement : aucune
évaluation ne trouve d'effet sur l'investissement des entreprises. Il se défend par ses recettes, par sa
lisibilité et parce qu'il réduit l'écart d'imposition entre un bénéfice distribué et un bénéfice conservé.

La faiblesse du système n'est pas dans le taux du PFU, mais à côté. Un bénéfice remonté dans une holding et
conservé y supporte environ 27~\% d'impôt au total, contre 49~\% s'il est distribué au ménage, et les
plus-values logées dans une holding s'effacent le plus souvent au décès. C'est là que porte la correction
proposée.

\section{Le PFU en 2026}

Le PFU additionne 12,8~\% d'impôt sur le revenu (art.~200~A du CGI) et 18,6~\% de prélèvements sociaux. Ce
second taux était de 17,2~\% jusqu'en 2025 : l'article 12 de la LFSS pour 2026 a porté la CSG sur les revenus
du capital de 9,2 à 10,6~\%, mais l'a laissée à 9,2~\% pour une liste fermée de revenus, dont l'assurance-vie,
l'épargne logement, les revenus fonciers et les plus-values immobilières. Il existe donc désormais deux taux
de prélèvements sociaux sur l'épargne, sans justification économique. Les loyers et les plus-values
immobilières restent hors PFU. Le contribuable peut opter pour le barème, globalement pour tous ses revenus
du capital de l'année ; l'option ouvre droit à l'abattement de 40~\% sur les dividendes et, depuis la LF
2026, elle est révocable. 2,5~millions de foyers l'ont exercée sur leurs revenus de 2024.

Deux contributions s'ajoutent en haut de la distribution. La contribution exceptionnelle sur les hauts
revenus (CEHR) prélève 3 puis 4~\% au-delà de 250\,000~€ de revenu fiscal pour une personne seule
(500\,000~€ pour un couple) : le taux marginal sur les revenus du capital monte alors à 35,4~\%. La
contribution différentielle (CDHR), créée pour les revenus de 2025, garantit que l'impôt sur le revenu et la
CEHR atteignent au moins 20~\% du revenu au-delà des mêmes seuils ; pour un foyer qui vit de dividendes,
elle porte le taux à 38,6~\%. Elle est prorogée jusqu'à ce que le déficit du budget de l'État repasse sous 3~\% du PIB, et
le PLF 2027 en attend 0,6~Md€.

Le PFU ne s'apprécie qu'avec l'impôt déjà payé par la société. Sur 100~€ de bénéfice, une société au taux
normal paie 25,8~€ d'impôt sur les sociétés et de contribution sociale ; si elle distribue le reste,
l'actionnaire paie 31,4~\% des 74,2~€ restants, soit 23,3~€. Le bénéfice distribué supporte donc 49,1~\% au
total, et jusqu'à 54,5~\% pour un foyer à la CDHR (figure~\ref{fig:taux}). C'est proche du taux marginal
supérieur sur un salaire, 55,4~\% hors cotisations sociales : pour un bénéfice distribué, l'imposition de la
société et celle de l'actionnaire s'emboîtent à peu près comme celle d'un revenu du travail. C'est élevé en
Europe. La Tax Foundation place la France au sixième rang sur 35 pour le taux personnel supérieur sur les
dividendes (34~\% en 2025, pour une moyenne de 21~\%). L'OCDE trouve, à cinq fois le salaire moyen, un taux
intégré de 49~\% que seuls le Danemark et l'Espagne dépassent (Hourani et al., 2023, cité par France
Stratégie).

Le chiffre qui surprend est au bas de la figure. Si la société verse son dividende à une holding détenue par
l'actionnaire et que celle-ci le garde, la holding ne paie que 1,25~\% : le régime mère-fille exonère le
dividende, sauf une quote-part de 5~\% imposée à 25~\%. Le bénéfice n'a alors supporté que 26,7~\%. Tant qu'il
reste dans la holding, il peut être placé ou réinvesti sans autre impôt.

\figbloc{pfu_01_taux_cumule}{Impôt total sur 100~€ de bénéfice, selon ce qu'en fait l'actionnaire (2026)}{fig:taux}

\section{Ce que les réformes de 2013 et de 2018 ont appris}

La France a changé deux fois de régime en cinq ans, en sens opposés, ce qui fournit une expérience rare. En
2013, l'option pour un prélèvement libératoire disparaît et tous les dividendes sont imposés au barème. L'IPP compare les foyers qui
avaient choisi le prélèvement libératoire en 2012, donc les plus touchés, aux autres foyers (Bach et al.,
2019). Leurs dividendes déclarés baissent de 40~\%, sans hausse de leurs salaires ni de leurs autres
revenus. Les entreprises qu'ils détiennent distribuent moins et accumulent des actifs financiers ; elles ne
réduisent pas leur investissement. La réforme devait rapporter 0,4~Md€ ; elle en a coûté 0,9.

Pourquoi le dividende réagit-il autant et l'investissement si peu ? Pour l'actionnaire d'une société qui
finance ses investissements sur ses bénéfices, l'impôt sur le dividende ne change pas le choix d'investir.
Prenons 100~€ de bénéfice après impôt sur les sociétés. Distribués aujourd'hui, ils rapportent 68,6~€ nets de
PFU. Réinvestis à 5~\%, ils donneront 105~€ distribuables l'an prochain, soit 72,0~€ nets. Le PFU prend la
même fraction dans les deux cas : il réduit le coût de l'investissement, le dividende auquel on renonce, dans
la même proportion que son rendement, et laisse le choix inchangé. Ce qui change, c'est la date et la forme du
revenu. Quand l'impôt sur le dividende monte, l'actionnaire qui contrôle sa société laisse le bénéfice
dedans : il le récupérera plus tard sous forme de plus-value, ou jamais s'il transmet ses titres. C'est
exactement ce qu'on observe en 2013. Le raisonnement ne vaut pas pour une jeune entreprise qui lève des fonds
propres auprès d'investisseurs extérieurs : l'impôt réduit alors le rendement sans réduire le coût. Le seul
effet réel trouvé après 2018 va dans ce sens : des créations d'entreprises un peu plus rapides dans les
secteurs les plus exposés, d'ampleur faible et difficile à séparer de la baisse de l'impôt sur les sociétés
(IPP, 2023). La baisse américaine de 2003 n'a pas davantage modifié l'investissement (Yagan, 2015).

Le PFU a produit l'effet inverse de 2013. Les dividendes reçus par les ménages passent d'environ 14~Md€ par an
en 2013-2017 à 23~Md€ en 2018-2020 ; l'IPP attribue à la réforme une hausse de 15~\% dès 2018. Le coût
budgétaire en est réduit d'autant : 0,4~Md€ au lieu de 0,9 à comportements inchangés selon l'IPP ;
0,74~Md€ au lieu de 1,05 dans l'hypothèse basse de l'Insee, et un gain de 0,36~Md€ dans son hypothèse haute
(Paquier et Sicsic, 2021). Le comité d'évaluation de France Stratégie conclut que le coût a été « annulé les
premières années » (2023). Cette conclusion repose sur la part de la hausse des dividendes qu'il attribue au
PFU, plus de la moitié, qui est un jugement et non une mesure. Le tableau~\ref{tab:etabli} résume ce que
les évaluations établissent.

\begin{table}[!htbp]\centering\small
\caption{Ce que les évaluations du PFU établissent}\label{tab:etabli}
\begin{tabularx}{\linewidth}{@{}>{\raggedright\arraybackslash}p{4.6cm}>{\raggedright\arraybackslash}X@{}}
\toprule
Effet & Ce qu'on sait \\
\midrule
Dividendes distribués & Forte réaction au taux : $-$40~\% en 2013, +15~\% en 2018. Établi par deux réformes de sens opposé \\
Coût budgétaire & Très inférieur au chiffrage statique ; autofinancement possible, non mesuré \\
Investissement des entreprises existantes & Aucun effet détecté, en 2013 comme en 2018 \\
Créations d'entreprises & Un peu plus rapides dans les secteurs exposés ; effet faible, mêlé à la baisse de l'IS \\
Salaires transformés en dividendes & Non observé au niveau agrégé \\
Exil & Faible dans le centile le plus riche ; non mesuré au-delà \\
\bottomrule
\end{tabularx}
\par\vspace{2pt}{\scriptsize\raggedright Sources : Bach et al. (2019) ; Paquier et Sicsic (2021) ; France
Stratégie (2023) ; IPP (2023) ; CAE, Focus n°~118 (2025).\par}
\end{table}

Le revers est la concentration. En 2021, 1~\% des foyers déclarent 96~\% des dividendes, 0,1~\% en déclarent
62~\% et 0,01~\% (environ 4\,000 foyers) en déclarent 33~\% (France Stratégie, 2023). Ces parts ont gagné 5 à
10 points depuis 2017, ce qui mesure d'abord un changement de ce qui est déclaré, des bénéfices autrefois
gardés dans les sociétés étant désormais distribués, plutôt qu'un enrichissement. Toute variation du taux du
PFU profite ou coûte donc presque entièrement à ce centile.

Reste à savoir si le taux est le bon. La seule estimation française du taux qui maximise les recettes sur les
revenus du capital, prélèvements sociaux compris, le situe entre 43~\% et 57~\%, selon qu'on tient compte ou
non de ce que les foyers déplacent aussi une partie de leurs revenus d'activité quand ce taux change
(Lefebvre, Lehmann et Sicsic, 2025). Le PFU, de 31,4 à 38,6~\%, est en dessous, mais la marge est mince.
Avec l'élasticité qu'ils estiment (0,77), un relèvement de 31,4 à 33~\% ne garde qu'un peu moins des deux
tiers de son rendement mécanique ; le CAE chiffre un relèvement de 3 points à 1,2~Md€ nets (Focus n°~124,
2025). L'exil est faible : 0,2~\% du centile le plus riche en revenus du capital part chaque année, et un
point d'imposition effective supplémentaire en fait partir à long terme 0,02 à 0,23~\% de plus (Bach et
al., Focus CAE n°~118, 2025). Aucune donnée ne mesure le comportement des 0,01~\% les plus riches. La
commission d'enquête de l'Assemblée sur l'imposition des plus hauts patrimoines recommande d'ailleurs de
« maintenir la stabilité » des taux du PFU (rapport n°~3056, recommandation~13).

\section{Deux idées à écarter}

\textbf{Ramener le PFU à 25~\%.} L'assiette du PFU est d'environ 80~Md€ (revenus de 2024 : 50~Md€ de
dividendes, 10~Md€ d'intérêts, 19~Md€ de plus-values). Baisser le taux de 6,4 points coûte 5,1~Md€ à
comportements inchangés. Avec une élasticité de 0,65 à 1, la hausse de l'assiette n'en rembourse qu'une
partie et le coût net reste de 3,2 à 3,9~Md€ par an. Ce gain irait pour l'essentiel au centile qui perçoit
96~\% des dividendes. Il ne corrige aucun défaut identifié : le frein aux distributions créé en 2013 a été levé
en 2018, et l'investissement n'a réagi ni à l'un ni à l'autre. Il creuserait en outre l'écart avec le
salaire (44~\% contre 55~\% sur la figure~\ref{fig:taux}) et pousserait les dirigeants-actionnaires à se
rémunérer en dividendes plutôt qu'en salaire. On n'observe pas ce comportement au niveau agrégé aujourd'hui,
mais il dépend précisément de cet écart.

\textbf{Un impôt de 25~\% sur les plus-values et les revenus non réinvestis des holdings.} Il vise le bon
problème, le circuit décrit au §~1, mais avec un instrument qui ne passe pas. Juridiquement, la directive
mère-fille (2011/96/UE) interdit d'imposer les dividendes qu'une filiale européenne verse à sa mère au-delà
d'une quote-part de 5~\%, que la France applique déjà ; et imposer chez l'associé des bénéfices qu'il n'a pas
perçus se heurte à la jurisprudence constitutionnelle, qui exige des revenus disponibles (analyse de la
direction de la législation fiscale rapportée par la commission d'enquête). Économiquement, il ajoute un
troisième étage d'imposition, après la société opérationnelle et avant l'associé, et oblige à définir ce
qu'est un revenu « non réinvesti ». Surtout, il alourdirait fortement l'impôt des très grandes fortunes, dont
85 à 90~\% détiennent une holding (IPP, note n°~92, 2023) : la quote-part de 1,25~\% deviendrait 25~\% sur une partie de
l'assiette. C'est l'inverse de la contrainte que se fixe ce rapport. Associé à la baisse du PFU, il déplacerait
la charge entre très riches, plus pour ceux qui gardent dans une holding, moins pour ceux qui se font
distribuer, sans gain d'efficacité démontré.

\section{Une correction et une articulation}

\subsection*{Faire expirer le report de l'apport-cession à la transmission}

Un entrepreneur qui vend sa société paie 31,4~\% sur sa plus-value. S'il apporte d'abord ses titres à une
holding qu'il contrôle, puis fait vendre la société par la holding, la plus-value d'apport est placée en
report : elle n'est imposée que si la holding vend sans réinvestir au moins 70~\% du produit dans une activité
économique (art.~150-0~B~ter du CGI, version de la LF 2026). Ce report se défend tant qu'il diffère l'impôt
sur un capital qui reste investi. Il cesse de se défendre à la transmission : au décès, le report devient une
exonération définitive ; en cas de donation, il passe au donataire et s'efface si celui-ci ne vend pas
pendant six ans. Les montants sont considérables. 33,6~Md€ de plus-values ont été mises en report en 2023 et
22,2~Md€ en 2024, davantage que les 18,8~Md€ de plus-values imposées au PFU. Le stock atteint 153~Md€, les
deux tiers des entrées de 2023 sont concentrés sur 84 foyers, et seulement 0,4 à 1,5~Md€ de reports expirent
chaque année (rapport n°~3056). Avec les bénéfices conservés dans les holdings, c'est ce qui fait baisser le
taux d'imposition total tout en haut de la distribution, selon la commission d'enquête.

L'article 5 du PLF 2027, en discussion au Parlement, fait expirer le report à toute transmission, à titre
gratuit comme onéreux, pour les transmissions intervenues depuis le 1\up{er} octobre 2026. Il permet d'étaler
le paiement sur cinq ans, pour que les héritiers n'aient pas à vendre l'entreprise pour payer, et fait expirer
de même le sursis de l'exit tax sur ces plus-values. La commission d'enquête le recommande (recommandation~11), et la direction
de la législation fiscale le juge sans obstacle constitutionnel, à la différence d'une imposition générale
des plus-values latentes au décès. Le rendement annoncé est d'environ 0,5~Md€ par an (chiffre de seconde
main, non vérifié dans l'évaluation préalable).

C'est la seule mesure de cette note qui alourdit nettement l'impôt de grandes fortunes : elle touche les
quelque 9\,000 foyers qui placent chaque année des plus-values en report, et une poignée d'entre eux très
fortement. Elle reste compatible avec la contrainte du rapport parce qu'elle ne relève aucun taux : elle
supprime un effacement, et un foyer qui garde sa holding jusqu'à sa succession paiera ce qu'aurait payé un
foyer qui vend, pas davantage. Elle a aussi un argument d'efficacité. Tant que la plus-value s'efface au
décès, un propriétaire âgé a intérêt à garder ses capitaux dans la holding jusqu'à sa succession plutôt qu'à
les distribuer ou à céder ses titres ; si le report expire de toute façon, attendre ne fait plus économiser
l'impôt. On peut supposer que cela libère des capitaux immobilisés, sans que l'effet ait été mesuré.

\subsection*{Garder la majoration de 4 points quand l'IR et la CSG fusionnent}

La note sur la fusion de l'impôt sur le revenu et de la CSG absorbe la CEHR dans le barème fusionné et reporte
sur les autres revenus la part payée sur les revenus du capital. Le PFU serait alors de 31,4~\% pour tous, au
lieu de 35,4~\% au-delà des seuils de la CEHR : une baisse de 4 points pour les foyers qui perçoivent
l'essentiel des dividendes, contraire à l'objectif de garder le PFU tel qu'il est. On garde donc cette
majoration sur les revenus du capital, aux mêmes seuils, la CEHR ne disparaissant que pour les revenus soumis
au barème fusionné. Celui-ci a un peu moins à lever : une partie, que les données publiques ne ventilent pas,
des 1,4~Md€ de CEHR. La CDHR reste en l'état jusqu'à son extinction prévue, mais son assiette (l'impôt sur le
revenu et la CEHR) devra être redéfinie dans l'impôt fusionné. On ne la durcit pas, comme le propose la
commission d'enquête (recommandations 14 à 17) : elle vise les mêmes foyers que l'expiration du report, et
cumuler les deux irait contre la contrainte du rapport.

\subsection*{Ce qu'on ne fait pas}

Trois mesures sont défendables mais ne sont pas retenues. Aligner les deux taux de prélèvements sociaux
supprimerait une distorsion sans fondement : un épargnant devrait choisir entre un compte-titres et une
assurance-vie selon leur rendement, pas selon leur fiscalité. L'alignement à 18,6~\% rapporterait de l'ordre
de 1,3~Md€ (estimation fragile), payés surtout par la classe moyenne épargnante ; c'est une recette
disponible si le rapport en a besoin, pas une mesure de croissance. Un prélèvement modéré sur la trésorerie
passive des holdings, distinguée de leur trésorerie opérationnelle, est la version défendable de l'impôt de
25~\% écarté plus haut ; l'Irlande taxe ainsi à 20~\% les revenus passifs non distribués. Il alourdirait lui
aussi l'impôt des grandes fortunes, et on le garde en réserve si l'expiration du report ne suffit pas. Enfin,
on ne rétablit pas d'abattement pour durée de détention sur les plus-values : celui de 2013-2017 avait
fortement compliqué l'impôt (France Stratégie, 2023). La correction de principe pour une détention longue est d'imposer
le gain réel, en indexant le prix d'acquisition sur l'inflation, ou d'exonérer le rendement normal du
capital comme le prévoit le cadre dual du rapport. Elle a un coût, non chiffré ici. Le PFU étant
proportionnel et sans quotient familial, rien dans cette note ne modifie la situation des familles.

\vspace{6pt}
{\scriptsize\textbf{Références.} Assemblée nationale (2026), commission d'enquête sur l'imposition des plus
hauts patrimoines et des revenus les plus élevés, rapport n°~3056 (J.-P. Mattei, président ; C. de Courson,
rapporteur). Bach, Bozio, Fabre, Guillouzouic, Leroy et Malgouyres (2019), note IPP n°~46. Bach, Bozio, Grimprel, Guillouzouic, Landais et Malgouyres
(2025), Focus du CAE n°~118. CAE (2025, 2026), Focus n°~124 et n°~147. France Stratégie (2023), Comité
d'évaluation des réformes de la fiscalité du capital, quatrième rapport. Hourani, Millar-Powell, Perret et
Ramm (2023), OCDE, \emph{Taxation Working Paper} n°~65. IPP (2023), rapport n°~47 sur le réinvestissement des
plus-values et les créations d'entreprises ; note n°~92, « Quels impôts les milliardaires paient-ils ? ». Lefebvre, Lehmann et Sicsic (2025), « Estimating the Laffer Tax
Rate on Capital Income: Cross-Base Responses Matter! », \emph{Scandinavian Journal of Economics} (IZA DP
16112, 2023). Paquier et Sicsic (2021), \emph{Économie et Statistique} n°~530-31. Tax Foundation (2026),
\emph{Dividend Tax Rates in Europe}. Yagan (2015), « Capital Tax Reform and the Real Economy: The Effects of
the 2003 Dividend Tax Cut », \emph{American Economic Review} 105(12). LFSS 2026 (loi n°~2025-1403), art.~12 ;
LF 2026 (loi n°~2026-103) ; PLF 2027, art.~5.\par}

\end{document}
